\documentclass[11pt]{article}
\usepackage[letterpaper,margin=1in]{geometry}
\usepackage{amsmath,amssymb,amsthm,booktabs}
\usepackage[hidelinks]{hyperref}
\setlength{\intextsep}{6pt plus 2pt minus 2pt}

\newtheorem{theorem}{Theorem}
\newtheorem{corollary}[theorem]{Corollary}
\newcommand{\tr}{\operatorname{Tr}}
\newcommand{\cF}{c_{\mathrm F}}
\newcommand{\pF}{\Pi_{\mathrm F}}
\newcommand{\CB}{\mathrm{CB}}

\title{Exact nonlocal stabilizer R\'enyi entropy at Schmidt rank seven}
\author{Heming Xiao\\[0.25em]
  \normalsize Independent Research}
\date{Revised 21 November 2026}

\begin{document}
\maketitle

\begin{abstract}
We give a computer-assisted proof that the computational-basis Schmidt
representative with descending Schmidt probabilities minimizes the
second-order stabilizer R\'enyi entropy over the complex local-unitary
orbit of every pure state of Schmidt rank seven. The result holds for
arbitrary qubit bipartitions, including those with unequal subsystem
dimensions, and gives an explicit formula for the minimum in terms of
the seven Schmidt probabilities. A paired-frame reduction and
dimension-independent Pauli estimates reduce the problem to interval
coverage conditions. We establish this coverage using outward-rounded
integer interval arithmetic in C++ and exact rational arithmetic in
Python. A formal proof in Lean~4.36.0 verifies the full interval
cover, the optimality of the sorted representative, and the formula
for the minimum.
\end{abstract}

\section{Problem and result}

Let \(A\) and \(B\) contain \(n_A\) and \(n_B\) qubits, respectively.
For a normalized pure state define its stabilizer purity and order-two
stabilizer R\'enyi entropy by
\[
 P_2(\psi)=2^{-(n_A+n_B)}
 \sum_{P\in\{I,X,Y,Z\}^{\otimes(n_A+n_B)}}
 |\langle\psi|P|\psi\rangle|^4,\qquad
 M_2(\psi)=-\log_2 P_2(\psi).
\]
The stabilizer R\'enyi entropy is a measure of nonstabilizerness
\cite{leone}. Its nonlocal version minimizes \(M_2\) over
\(U_A\otimes U_B\) \cite{cao,qw,torre,lcold}. Let
\(t_0\ge\cdots\ge t_6>0\) be the nonzero Schmidt probabilities and let
\(\lvert\psi_{\CB}\rangle=\sum_{j=0}^{6}\sqrt{t_j}\lvert j\rangle_A
\lvert j\rangle_B\), using binary labels with leading zeros as needed.

The ordered Schmidt gauge is a proposed minimizer
\cite{torre,lcold,fg,sierant}. Sierant proved the conjecture for all
spectra of rank at most six and for dyadic-staircase spectra
\cite[Theorems 1 and 2]{sierant}. Section II.2.2 of that work
identifies general rank seven as an open case. Liu and Cui proved local
minimality in the ordered gauge for integer orders at least two
\cite[Eq.~(5) and Appendix C]{lcnew}. The local-orbit optimization
based on stabilizer fidelity also has an exact solution
\cite{sfid,viscardi}. Huang, Chen,
and Yao study dimension reduction and compute the rank-two value
\cite{hcy}. The result here uses a comparison on every ambient qubit
dimension.

\begin{theorem}\label{thm:rank7}
For every pure state of Schmidt rank seven on any qubit bipartition,
\[
 \min_{U_A,U_B} M_2((U_A\otimes U_B)\psi)=M_2(\psi_{\CB}),
\]
where the minimum includes all complex local unitaries.
\end{theorem}

\begin{corollary}\label{cor:formula}
Let \(\mathcal A\) be the four-element subsets of
\(\{0,\ldots,6\}\) whose binary labels have XOR zero. The minimum in
Theorem~\ref{thm:rank7} equals
\[
 -\log_2\!\left[
 7\left(\sum_{j=0}^{6}t_j^2\right)^2
 -6\sum_{j=0}^{6}t_j^4
 +168\sum_{I\in\mathcal A}\prod_{i\in I}t_i
 \right].
\]
\end{corollary}

\section{Reduction to a Pauli-defect bound}

On a balanced cut of local dimension \(D\), write the coefficient
matrix as \(T=U_A M U_B^{\mathsf T}\), where
\(M=\operatorname{diag}(\sqrt{t_0},\ldots,\sqrt{t_6},0,\ldots,0)\).
Set \(H=M^{\otimes4}\) and \(J=D^{-2}\sum_P P^{\otimes4}\), with
the sum over the sign-free single-side Pauli strings. The purity is
\[
 D^2\tr(HAHB),\quad
 A=(U_A^\dagger)^{\otimes4} J U_A^{\otimes4},\qquad
 B=(U_B^{\mathsf T})^{\otimes4} J \overline{U_B}^{\otimes4}.
\]
This follows by expanding
\(\langle P\otimes Q\rangle=\tr(T^\dagger P T Q^{\mathsf T})\)
to the fourth power, using \(J^{\mathsf T}=J\).
Hilbert--Schmidt Cauchy--Schwarz bounds the expression by
\(\sqrt{D^2\tr(HAHA)\,D^2\tr(HBHB)}\).
The two factors are actual paired-frame purities. A paired frame has
coefficient matrix \(N=VMV^\dagger\) for a unitary \(V\).
The paired reduction is also given in
\cite[Lemma 1 and Corollary 1]{sierant} and
\cite[Eq.~(4) and Appendix B]{lcnew}.
For unequal qubit sides, appending leading computational ancillas to
the smaller side embeds the states being compared in a balanced
orbit. Each ancilla contributes a factor of one to \(P_2\), and the
sorted representative uses the same binary labels after padding
\cite[Eq.~(16)]{sierant}.

Let \(\pF\) be the antisymmetric Pauli projector of
\cite[Eq.~(39)]{sierant}. Define
\(\cF(K)=\tr(\pF\Lambda^4K)\), and put \(\eta=4/D\). Choose the
orthonormal eigenframe \(u_0,\ldots,u_6\) for the positive spectrum
of \(N^2\). For a four-element set \(I\), define
\[
 w_I=\eta^{-1}\langle u_I,\pF u_I\rangle,\qquad
 E=\sum_{|I|=4}(w_I-w_I^{\CB})\prod_{i\in I}t_i,
\]
where \(u_I\) is the exterior product of the \(u_i\) in index order.
The computational array \(w_I^{\CB}\) is the indicator of
\(\mathcal A\), so \(\cF(N^2)-\cF(M^2)=\eta E\).
Sierant's sector decomposition and rank-seven compression bound give
\begin{align*}
 G(V)&=\mathcal I(t)+6D^2 Q_{(1^4)}(N)
                  +18D\,\cF(N^2)-\mathcal D(N),\\
 \mathcal I(t)&=7\left(\sum_jt_j^2\right)^2-6\sum_jt_j^4,
 \qquad Q_{(1^4)}(N)\le\eta\,\cF(N^2).
\end{align*}
Here \(\mathcal D\) is the nonnegative sum of the three sector
defects, which vanishes at \(V=I\), and the compression bound is
saturated there
\cite[Proposition 2, Eq.~(175), and Theorem 7]{sierant}.
Retain the normalized compression loss
\[
 \mathcal D_F(N)
 =\frac{\eta\cF(N^2)-Q_{(1^4)}(N)}{\eta^2}\ge0 .
\]
Then the exact purity difference and the target bound are
\begin{equation}\label{eq:defect}
 G(V)-G(I)=168E-\big(\mathcal D(N)+96\mathcal D_F(N)\big)\le0.
\end{equation}
The coefficients are \((6D^2\eta+18D)\eta=168\) and
\(6D^2\eta^2=96\). The sum of the two defects is the physical
quantity bounded by the certificate. Evaluating the
decomposition at \(V=I\) gives the expression inside the logarithm in
Corollary~\ref{cor:formula}. Its interpretation as the maximum follows
from Theorem~\ref{thm:rank7}.

\section{The finite certificate}

The accompanying proof supplement
(\texttt{rank7\_submission\_proof\_guide.md}) indexes the analytic
comparisons and verifier sources.
The following describes the inequalities certified by those sources and
how they imply Eq.~\eqref{eq:defect}.

The normalized four-wedge compression \(\mathsf C\) is a
positive contraction on the exterior space of the seven
Schmidt vectors \cite[Theorem 22]{sierant}. Its diagonal is
\((w_I)\). With \(a_I=\sqrt{\prod_{i\in I}t_i}\), the exact loss is
\[
 \mathcal D_F(N)=
 \sum_I a_I^2(\mathsf C-\mathsf C^2)_{II}
 +\sum_{I<J}(a_I-a_J)^2|\mathsf C_{IJ}|^2.
\]
For a coordinate cut with projection \(R\), the cut trace loss
is bounded below by \(\{s_R\}(1-\{s_R\})\), where
\(s_R=\tr(R\mathsf C R)\) and braces denote fractional part.
Two nested cuts can use disjoint amplitude intervals, sharing
the displayed compression loss. The shared cut requires an
explicit correction for exterior labels that cross its nominal
endpoint. The separate Pauli-sector estimates control the signed
mixed terms and the deficit left after these interval payments.

\paragraph{Spectral reduction.}
Write the ordered probabilities as \((\alpha,\beta,\gamma,a,b,c,d)\),
and let \(h=\alpha-\beta\), \(k=\beta-\gamma\), \(q=\gamma-a\),
\(p=a-b\). Normalize \(b=1\) or \(d=1\) as specified in the certificate.
All comparisons are homogeneous of degree four in the probabilities.
The two defects are nonnegative, so \(E\le0\) immediately gives
Eq.~\eqref{eq:defect}. On \(E>0\), the mixed branch certifies the
stronger bound with \(A_0E\) in place of \(168E\), where
\(A_0=(101/100)168\). The weighted and analytic pruning branches
establish Eq.~\eqref{eq:defect} directly. The analytic criteria
cover either nonpositive leading trace excess and specified spectral regions.
In particular, the unconditional small-gap comparison handles
\(q\le\max(2d/3,b/2)\). The remaining variables
\((h,k,q,c,d)\) lie inside explicit bounded boxes after rescaling.

\paragraph{Transport from \(a=b\).}
For a fixed physical frame, let \(F=E-\lambda U-\mu U_1\) be
the signed exterior gain after the two mixed terms, with
\(\lambda=24s/A_0,\ \mu=24t/A_0,\ A_0=(101/100)168\).
The mixed comparison at \(p=0\) gives
\[
 A_0F(0)+96D_{\mathrm{seed}}
 \le Q+96\big(W_1(0)T_1+W_s(0)T_s\big).
\]
Here \(Q\) is a physical Pauli remainder usable throughout the ray,
\(W_1,W_s\) are disjoint-interval squared widths, and \(T_1,T_s\)
are lower trace prices bounded above by \(e=1-w_{0123}\).
The seed allowance \(D_{\mathrm{seed}}\) covers the shared
crossing at \(p=0\). The decay of the widths and the increase in
crossing error are bounded by the decrease of \(A_0F(p)\) for
\(q\ge2d/3\), with \(0\le s\le3hq/2\) and
\(0\le t\le7k(h+k+q)/4\). Hence the same remainder and trace
prices certify Eq.~\eqref{eq:defect} for every \(p\ge0\).
The independent weighted branch has its own transfer bound,
including the possible loss of its highest width. The older
pruning criteria either hold at every \(p\) or transmit by spectral
monotonicity. These three cases are identified separately in the
supplement for both implementations.

\paragraph{Interval comparison on the tie.}
In the upper-ratio verifier \(a=b=1\), the underlying signed
exterior gain plus its crossing allowance is multiaffine in the
two tail probabilities \(c,d\) on the same physical frame.
The lower bound on the defect uses a common multiaffine minorant.
At the four ordered corners, each of the two exterior estimates
can be mixed with the other using a corner-specific weight.
The same shifts work at all four corners and on both boundary
rays of each sector of \(v/n\), where \(v,n>0\) are the two trace
excesses. The ray boundaries are \(0,1/2,2/3,1,2,4,\infty\).
Concavity of the trace lower prices on each bounded sector extends
the inequalities between its endpoints. The unbounded sector
uses an affine comparison and a monotone trace lower price.
The bounded-corner lower price and the actual signed gain then
interpolate throughout the ordered rectangle.
For the lower-ratio verifier only the sixth weight is interpolated,
and the highest compression price may be set to zero.

\paragraph{Exhaustive cover.}
Both verifiers either accept a box by a sufficient comparison,
remove a region covered by the analytic reductions, or bisect
an interval into two closed subintervals. Every accepted comparison
uses outward-rounded integer interval endpoints. No sampled margin
is needed. The proof of each contractor states the necessary
inequality it preserves. Table~\ref{tab:cover} gives the exact
ratio intervals, where \(\xi=d/b\). The two end regions use the
large-ratio and close-bottom results proved in the supplement.

\begin{table}[h]
\centering
\begin{tabular}{@{}lll@{}}
\toprule
Range of \(\xi\) & Comparison & Recorded leaves \\
\midrule
\((0,1/2]\) & large fifth-to-seventh ratio & \(93{,}703\) \\
\([1/2,9/16]\) & upper verifier, first slab & \(6{,}455{,}151\) \\
\([9/16,3/4]\) & upper verifier, thirteen closed cells & \(127{,}384{,}781\) \\
\([3/4,100/101]\) & lower verifier, six closed slabs & \(4{,}350{,}568\) \\
\([100/101,1]\) & close bottom probabilities & see supplement \\
\bottomrule
\end{tabular}
\caption{Ratio intervals in the external verification. The first four
rows count accepted leaves, including boxes pruned by proved criteria.}
\label{tab:cover}
\end{table}

\begin{proof}[Proof of Theorem~\ref{thm:rank7}]
Table~\ref{tab:cover} covers every \(0<\xi\le1\), including all
joining endpoints. The transport argument covers arbitrary \(a-b\)
where a tie certificate is used. Therefore the Pauli-defect bound
\eqref{eq:defect} holds for every paired frame and every spectrum of
rank seven. Cauchy--Schwarz transfers the resulting upper purity bound
to the full local-unitary orbit. Leading computational padding gives
the same result on unequal cuts. The sorted representative attains
the bound, so it attains the minimum of \(M_2\).
\end{proof}

\section{Reproducibility and disclosure}

The \href{https://github.com/heming-xiao/rank-seven-stabilizer-entropy/blob/a9daab9042819dfbe8b26dd4926ffd9a9079dad9/proof/rank7_submission_proof_guide.md}{proof guide for the external verifiers}
lists the analytic lemmas, source files, run records, and replay
commands. The coverage checker joins the thirteen closed intervals
using rational arithmetic and checks the source hashes and run
summaries. Replaying the C++ verifier regenerates the interval
traversal. Separate AI-assisted reviews examined the analytic
reductions, the verifier code, and the interval records.

The Lean~4.36.0 proof in the v5 source bundle starts with a normalized
state given by a seven-term Schmidt decomposition. It allows arbitrary
qubit dimensions on both sides and complex Schmidt frames. The Lean
kernel checks the full interval cover and the deductions from the
Pauli estimates to Theorem~\ref{thm:rank7}, including attainment by
the sorted representative and the formula in Corollary~\ref{cor:formula}.
The external verifier counts in Table~\ref{tab:cover} describe a
separate calculation.

In the v5 sources, \texttt{localUnitaryRepresentative} proves
Theorem~\ref{thm:rank7} and \texttt{localUnitaryMinimum} gives the
formula in Corollary~\ref{cor:formula}. All 5,683 project modules
passed strict compilation from source. The audit lists only
\texttt{propext}, \texttt{Classical.choice}, and \texttt{Quot.sound}
across seven declarations. The reproduction bundle includes the pinned
dependency versions, build instructions, and the full output in
\path{audit/Rank7TheoremAudit.log}. The manuscript and the complete v5
proof bundle are archived at
\href{https://doi.org/10.5281/zenodo.24421820}{\texttt{doi:10.5281/zenodo.24421820}}.

\newpage
\paragraph{Generative AI disclosure and author responsibility.}
Heming Xiao initiated and directed the research project. OpenAI Codex
agents using GPT-6.2 Sol and GPT-6.2 Astra developed the mathematical
arguments and proof, wrote the verification code, and drafted the
manuscript.

The proof has undergone AI-assisted review and computational checks.
Independent verification of the complete proof by the author and by
external experts remains pending. Heming Xiao takes responsibility for
the mathematical content of the manuscript and the accuracy of the AI
disclosure, and will publicly correct any errors identified in the
manuscript.

\begin{thebibliography}{10}
\bibitem{leone}
L.~Leone, S.~F.~E. Oliviero, and A.~Hamma,
Stabilizer R\'enyi Entropy,
\href{https://doi.org/10.1103/PhysRevLett.128.050402}{Phys. Rev. Lett. \textbf{128}, 050402 (2022)}.

\bibitem{cao}
C.~Cao, G.~Cheng, A.~Hamma, L.~Leone, W.~Munizzi, and S.~F.~E. Oliviero,
Gravitational Backreaction Is Magical,
\href{https://doi.org/10.1103/z3vr-w5c5}{PRX Quantum \textbf{6}, 040375 (2025)}.

\bibitem{qw}
D.~Qian and J.~Wang,
Quantum nonlocal nonstabilizerness,
\href{https://doi.org/10.1103/PhysRevA.111.052443}{Phys. Rev. A \textbf{111}, 052443 (2025)}.

\bibitem{torre}
G.~Torre, F.~Franchini, and S.~M.~Giampaolo,
Non-Local Magic from the Entanglement Spectrum,
\href{https://arxiv.org/abs/2607.07808v2}{arXiv:2607.07808v2} (2026).

\bibitem{lcold}
L.-Y.-N.~Liu and J.~Cui,
Nonlocal nonstabilizerness for slightly entangled quantum many-body states,
\href{https://arxiv.org/abs/2607.10714}{arXiv:2607.10714} (2026).

\bibitem{fg}
F.~Franchini and S.~M.~Giampaolo,
Schmidt-gauge non-local magic: Representation, optimality, and mathematical properties,
\href{https://arxiv.org/abs/2608.02745v2}{arXiv:2608.02745v2} (2026).

\bibitem{sierant}
P.~Sierant,
Spectral geometry of nonlocal stabilizer entropy,
\href{https://arxiv.org/abs/2609.13106v2}{arXiv:2609.13106v2} (2026).

\bibitem{lcnew}
L.-Y.-N.~Liu and J.~Cui,
Entanglement, anti-flatness, and nonlocal nonstabilizerness:
a unified perspective from entanglement spectrum,
\href{https://arxiv.org/abs/2609.01993v2}{arXiv:2609.01993v2} (2026).

\bibitem{sfid}
P.~Sierant,
Exact quantification of nonlocal magic,
\href{https://arxiv.org/abs/2608.28563}{arXiv:2608.28563} (2026).

\bibitem{viscardi}
M.~Viscardi, L.~Leone, and A.~Hamma,
Non-local Magic: closed-form solution and equivalence with magic of purification,
\href{https://arxiv.org/abs/2609.04119v2}{arXiv:2609.04119v2} (2026).

\bibitem{hcy}
X.~Huang, G.~Chen, and Y.~Yao,
Intrinsic spectral structure of bipartite nonlocal magic resource,
\href{https://arxiv.org/abs/2606.24368v1}{arXiv:2606.24368v1} (2026).
\end{thebibliography}
\end{document}
